\section{Discussion}
\label{ch:discussion}

The controlled study provides evidence about assistance on selected institutional questions, not a validated measure of deployed chatbot success. Its strongest result is that substantially increasing document context did not produce better AI-assessed answering or guidance; the absolute totals require more caution. This chapter answers the research questions, states the threats to validity and identifies the next evaluation.

\subsection{RQ1: observed assistance and remaining shortcomings}

Under G0, the prespecified combined endpoint assigns assistance to 31 of 96 email responses and 53 of 96 production-question responses. These totals describe retained AI labels, not the proportion of student questions the assistant factually answered. The email total comprises nine supported-answer and 22 useful-routing labels; the production total comprises 25 supported-answer and 28 useful-routing labels. Appropriate guidance therefore accounts for a substantial part of the positive result. A referral or clarification can be the right response when a question depends on missing information, an individual record or an institutional decision, but it does not resolve the underlying problem, which is why answering and routing are reported separately.

The remaining labels indicate concrete weaknesses. Under G0, 31 email and 22 production outcomes are incomplete, and the judge assigns the unsupported-claim category to twenty email and twelve production outcomes out of 96 each. Completed generation is therefore an insufficient success criterion, although the counts identify kinds of failure rather than their true prevalence.

The supported-answer category also includes responses without substantive factual targets: twenty of the 25 production supported-answer labels under G0 use references with no retained facts. A greeting or useful request for context may deserve a positive assessment, but should not count toward factual accuracy; an empty reference can also reflect failed reference construction. RQ1 is therefore answered descriptively: the system provides both answering and guidance, alongside substantial incompleteness and unsupported-claim labels, while the factual validity of the absolute rates remains unresolved.

\subsection{RQ2: the effect of retrieval policy}

Moving from G0 to S0 lowers the document threshold while holding the FAQ threshold and other selection rules fixed. Assistance changes from 31 to 30 email outcomes and from 53 to 52 production outcomes, a paired difference of $-1.04$ percentage points in each sample. The 95\% intervals, $[-8.33,6.25]$ and $[-7.29,5.21]$, span zero and allow improvement or deterioration. The evidence shows no benefit from this intervention, but neither does it show equivalence or identify a universally preferable threshold.

The small outcome differences occur despite pronounced changes to the retrieved input. For email questions, mean context length increases from 5,110 to 16,481 characters and empty contexts fall from 36 to three. For production questions, mean length increases from 1,010 to 5,808 characters and empty contexts fall from 76 to 52. In the email subset with nonempty checked references, reference-quotation availability also increases. The intervention thus changes the evidence presented to generation without a corresponding improvement in assistance.

Added passages may be irrelevant or redundant, the remaining difficulty may concern clarification or individual applicability, or the reference may omit a required condition; the experiment does not isolate these mechanisms. The result supports examining how answers use evidence rather than treating context volume as an effectiveness measure.

Fallback has to be interpreted through actual exposure. With the shared threshold, G1 changes empty contexts to nonempty ones in 33 email and 27 production cases, yet assistance changes by only $+1/96$ and $-1/96$. With the lower document threshold, S0 and S1 have identical contexts for every email case and for 93 of 96 production cases, so the absence of an email difference follows from identical prompts and answer reuse and says nothing general about fallback. Candidate allocation, neighboring passages, ordering and citation visibility remain untested.

The prompt comparison points the same way from the generation side. A prompt that checks source dates and applicability, preserves procedural steps and states gaps changes strict completion from 6/40 to 5/40, a $-2.5$-point difference with an interval spanning zero (Section~\ref{sec:results-prompt}). The comparison also yields a methodological lesson: under initial labels produced without the full generation evidence, the same answers appeared to show a five-point gain, so an assessment missing the generator's evidence can itself manufacture an improvement.

\subsection{RQ3: dependence on the sample and assessor}

The email and production samples should be read as two question frames. Their equal denominators do not represent the frequency of question types in actual use. Email questions come from correspondence groups already examined in the component study; production questions come from selective text logging, and the replay keeps only the latest turn, so a short follow-up may lack context from the original conversation. Differences between the samples cannot be assigned solely to language, script, retrieval quality or channel.

The assessment procedure adds a second source of uncertainty. Luna constructs references, generates answers and supplies primary judgments. Reference creation precedes the evaluated answers and policy labels are hidden from the judge, which reduces direct policy cues, but the shared model family can leave common omissions and interpretations across roles. Research on self-preference motivates examining this possibility, but does not identify its presence or magnitude here~\cite{panickssery2024selfpreference,chen2025selfpreference}.

Repeated Luna judgments agree on the overall category in 35 of 47 eligible pairs; local Qwen judgments agree with the primary categories in thirteen of 45. One assessment is therefore not perfectly stable, and the same saved evidence supports substantially different automated judgments. Neither comparison supplies human truth: Qwen's lower agreement does not justify replacing the primary labels, and Qwen applying the unsupported-claim category more broadly does not validate Luna.

RQ3 is thus bounded: the study gives interpretable conditional comparisons but weak support for precise absolute quality estimates. ARES, for example, uses human-labeled validation data within its evaluation framework~\cite{saadfalcon2024ares}; the present study has no corresponding human validation of its main labels.

\subsection{Construct validity and reference authority}

The prespecified endpoint deliberately combines supported answering with useful safe routing. This matches a broad assistance objective but is not a factual-accuracy measure, and it is sensitive to empty factual targets. The post-hoc sensitivity analysis requires a nonempty valid reference for the direct-answer category while retaining routing; production assistance then changes from 33/96 under G0 to 32/96 under S0. This alternative excludes some legitimate conversational responses and is not corrected accuracy.

The implemented judge returns an overall category, coverage and support flags, a rationale and evidence quotations, not the exhaustive claim-by-claim checks envisaged in the initial design, so a coverage flag is the assessor's assertion rather than a verified inventory of required facts. Structural and quotation checks establish well-formedness and source occurrence, not semantic entailment, exhaustive coverage or current applicability.

Authority is a separate issue from textual fidelity. Faculty transfer provisions from 2013 and university provisions from 2019 set different credit requirements for overlapping transfer situations, and a 2022 supplement permits an exception for students transferring from a country under a declared state of emergency, crisis or war~\cite{finki2013study,ukim2019study,ukim2022supplement}. Their presence in a retrieval pool does not decide which applies to a particular student. The fee schedule specifies applicability from 2024/25 until replacement~\cite{finki2024fees}; an upload date cannot substitute for that effective scope. The architecture must preserve source identity, dates and authority so they can be reviewed; the study does not certify current rules or individual eligibility.

The offline reference audit makes these weaknesses inspectable without retrospectively repairing the labels. Of 49 rejected main-reference quotations, 32 appear in another block in the same evidence pool and seventeen are absent.

\subsection{Internal, statistical and external validity}

Internal comparison is strengthened by paired questions, frozen GPU results, captured prompts and reuse of identical generated answers, with input rejections and unusable judgments kept in policy denominators. These controls do not remove reference errors shared by all policies. Separate read-only batches do not guarantee one historical database snapshot, and the intervention is specific to the selected corpus, candidate pool and generation configuration. The main experiment has no retrieval-disabled baseline, so it compares selection policies without measuring the contribution of retrieval itself. The cross-model check used a smaller local open-weight model; a comparably capable independent assessor would be more informative.

Statistical uncertainty is conditional on the saved AI judgments. Bootstrapping question groups preserves within-question pairing but does not propagate reference-construction uncertainty or systematic assessor error, and identical prompts are not independent answer draws merely because several policy records point to them. The intervals support a bounded inference about the observed intervention; their inclusion of zero is not a claim of equivalence.

External validity is restricted by both sampling and execution scope. Previously examined email questions are not a new holdout, and selectively recorded production text is not a random sample of interactions. The experiment replays production retrieval and prompts with controlled hosted generation, excluding live tool calls, conversation history, admission checks, storage and client delivery, so it does not cover full conversations, authenticated actions or the deployed agent. A shared reference architecture identifies plausible reuse points but, without another evaluated implementation, does not demonstrate portability to a second institution.

Historical telemetry documents delivery events and timing, but incomplete request linkage, selective feedback and its timing boundary prevent it from validating answer quality. Administrative workload, avoided emails and staff-time savings are outside this thesis's effectiveness scope.

\subsection{Attributing observed failures to architectural stages}

The architecture is most useful where it makes failure sources distinguishable: source preparation determines what can be supported, retrieval which evidence reaches generation, the answer stage how conditions are expressed, and channel adapters what the user sees. Mapping the observations onto these stages and onto the failure points that Barnett et al. report for RAG systems~\cite{barnett2024seven} shows where the remaining weaknesses lie. Table~\ref{tab:failure-attribution} summarizes the mapping; it is an analytical attribution from aggregate evidence, not a case-by-case causal diagnosis.

\begin{table}[htbp]
\centering
\caption{Observed failures attributed to architectural stages. FP numbers follow Barnett et al.~\cite{barnett2024seven}; rows without an FP number are additions of this study. Evidence refers to the controlled experiment unless stated otherwise.}
\label{tab:failure-attribution}
\small
\begin{tabularx}{\linewidth}{@{}>{\raggedright\arraybackslash}p{0.17\linewidth}>{\raggedright\arraybackslash}p{0.17\linewidth}>{\raggedright\arraybackslash}X@{}}
\toprule
Stage & Failure point & Observed evidence \\
\midrule
Source preparation & Missing content (FP1) & Under S0, 52 production contexts remain empty, although some are conversational turns that need no source; seventeen of 49 rejected reference quotations are absent from the evidence pool; overlapping transfer rules and dated fee schedules require authority and currentness metadata. \\
Evidence selection & Missed top-ranked; not in context (FP2, FP3) & Under G0, 36 email contexts are empty and on average fewer than half of the reference quotations reach the context; the separate document threshold raises this to about two thirds and improves ranked retrieval in the component study. \\
Answer construction & Not extracted; incomplete (FP4, FP7) & With far more evidence under S0, email incompleteness changes only from 31 to 30 and unsupported-claim categories from twenty to nineteen; the revised prompt does not improve audited completion. \\
Channel delivery & Input boundary; presentation (cf.\ FP5) & Four email questions exceed the 2,000-character input limit; the Discord adapter ignores the separate source event. Delivery behavior was not experimentally evaluated. \\
Observation and assessment & Evaluation evidence & Empty references, judge inputs lacking generation evidence and low cross-assessor agreement limit what the labels establish. \\
\bottomrule
\end{tabularx}
\end{table}

The mapping points to a specific engineering conclusion. The separate threshold largely resolves the evidence-selection failure for email questions: empty contexts nearly disappear and more reference quotations reach the model. Assistance nevertheless stays almost unchanged, which locates the dominant remaining weakness in answer construction, together with missing or ambiguous source content. For this system, further loosening of selection is unlikely to be the most productive change. Completeness checks during answer construction, explicit clarification when conditions are missing, and authority and currentness metadata at source preparation are better candidates. These are design implications derived from the stage separation, not interventions shown here to improve the pipeline; preserving the stage boundaries is what makes such attribution possible.

\subsection{Next evaluation}

The first proposed extension is to review references before generating new answers. Each case should receive an explicit factual, conversational, clarification, personal-action, mixed or out-of-scope task type. Factual and mixed cases need required facts with checked sources and applicability; other cases need explicit expected behavior, while any factual assertions still require support. Reviewer identity, AI assistance, human decisions and adjudication should remain distinct. Above all, qualified human review of a subset of answers, compared against the AI labels, would provide the currently absent evidence about label reliability; it is the most important next step for this line of work.

A subsequent comparison should freeze the baseline, minimum effect of interest, sample and analysis before execution. A retrieval-disabled condition could clarify the contribution of retrieved text, provided that useful links and agent tools are explicitly controlled. A stronger completeness intervention than the tested prompt could combine explicit identification of missing conditions, clarification requests and checks of displayed citations, assessed with the full generation evidence. The post-evaluation retrieval changes (Section~\ref{sec:post-evaluation-changes}) should be evaluated with the same controlled design. Because the existing cases and outcomes have been inspected, their reuse would be exploratory; a new held-out claim requires a newly selected and frozen sample.

Finally, a separate integration set should examine complete conversations and read-only tool tasks, using dated captured tool records to check tool selection, arguments, returned evidence and final answers, and keeping browser and Discord differences in history and source presentation explicit. Such work would extend the evidence beyond the controlled retrieval study reported here.
